\section{引言}
\label{sec:intro-zh}

中医远程舌诊需要机械臂载相机可重复地采集高质量舌象，方可支撑云端分析与电子病历（EMR）集成~\cite{zhang2020tongue,qi2016tongue}。
临床舌诊依赖颜色、纹理、苔质与形态等特征，在微动模糊、光照不均与拍摄距离失当时显著退化~\cite{zhang2020tongue}。
朴素流水线将每帧上传云端，既浪费带宽，又在链路较慢时放大端到端时延。

\paragraph{动机。}
云边协同机器人可将重型视觉模型卸载至云端，同时保持采集闭环响应性。
但在中医成像中 \emph{GIGO} 效应尤为突出：一次模糊上传即浪费 RTT，并可能将低置信特征传入下游分类器。
边侧质量门控因此须满足：（i）足够快以嵌入闭环；（ii）可解释、可审计；（iii）能在真实舌象语料上标定。
ShezhenV3-COCO 提供舌体边界框，使 ROI 裁剪与下游解剖关注区对齐。

\paragraph{方法概述。}
本文研究 FR3 平台上的\textbf{云--边--端}架构，采用 ROS~2 与 MoveIt。
轻量 \emph{Edge-IQA} 在 COCO 舌区 ROI 上融合 Laplacian 清晰度与曝光；LangGraph 在预算 $K$ 内决策三分支。
可选基线 B3 在训练集上微调 MobileNetV3-Small 作对比。

\paragraph{全量数据集评估。}
在 ShezhenV3-COCO（6{,}719 张）上：572 对验证样本（ROI + 合成模糊）标定 $\tau$；553 张测试图、每 seed 500 帧，使用真实 Edge-IQA 与 B3 分数。

\paragraph{贡献。}
\begin{enumerate}
  \item \textbf{Edge-IQA + ROI：}在舌区 crop 上标定 $\tau^\ast{=}0.465$，CPU p95 $<$10\,ms。
  \item \textbf{LangGraph 路由：}B2 使 M2$=1.00$，相对 B0（M2$=0.56$）显著提升。
  \item \textbf{真实 B3 对比：}MobileNet 验证准确率 98.3\%；共享 $\tau$ 下 M2$=0.829$，说明 $\tau$ 与可解释分数标度耦合。
  \item \textbf{双轨评估：}ShezhenV3 主矩阵 + Franka M6 辅轨，CSV 不混用。
\end{enumerate}

\paragraph{章节安排。}
第~\ref{sec:related-zh}--\ref{sec:conclusion-zh} 节分别综述、方法（含 ROI/B3）、实验与结论。

\paragraph{范围与非目标。}
本文不提出新舌病分类器、不宣称通用 IQA 排行榜 SOTA、不替代临床判断；聚焦采集时刻的质量门控。

\paragraph{可复现性。}
\texttt{paper1\_run\_tcm\_full.sh} 一键重建图表与 PDF；工件见 \texttt{experiments/results/}。

\paragraph{临床转化路径。}
试点应记录 50+ 会话 JSONL、本地重标定 $\tau$，并对比路由前后云端分类置信度。

\paragraph{与远程医疗工作流关系。}
人工重拍类似 \texttt{resample\_edge}，但缺乏统一 $Q_{\mathrm{img}}$ 阈值；Edge-IQA 标准化决策并导出 \texttt{flags}。

\paragraph{工程维护。}
发布说明中固定 \texttt{scorer.py} 与 checkpoint 的 SHA256；修改 ROI padding 须重跑标定。

\paragraph{教学价值。}
双轨设计分离算法证据（ShezhenV3）与集成证据（Franka M6），可复用于其他机器人传感论文。

\paragraph{数据与代码可用性。}
划分与脚本在仓库；ShezhenV3 图像路径由 \texttt{tcm\_paths.yaml} 配置。

\paragraph{致谢占位。}
相机就绪版补充基金与合作 clinic。

\paragraph{Closing.}
Paper~I 主张在真实 TCM  imagery 上验证的\emph{路由架构}，而非新的 IQA 回归 SOTA。
